\section{Subcritical uniformity and a Pell counterexample gate}
\label{sec:dual-subcritical-pell}

This section runs the positive and counterexample programs against the same
unchanged target.  The positive result identifies the exact pointwise
uniformity statement equivalent to $abc$.  The counterexample analysis shows
why fixed residual coefficients and polynomial identities cannot supply an
unbounded strict mixed-full family, while retaining one explicit conditional
Pell route.  No condition introduced below is asserted to hold without proof.

\subsection{The exact positive gate}

For an $abc$ point $P$, retain
\[
             H(P)=\log c,\qquad R(P)=\log\rad(abc).
\]
Consider the pointwise uniformity statement
\begin{equation}\label{eq:subcritical-locus-property}
 \forall\mu\in[0,1)\quad\exists B_\mu\in\R\quad\forall P,
 \qquad R(P)\leq\mu H(P)\Longrightarrow H(P)\leq B_\mu.
\end{equation}

\begin{theorem}[Subcritical-locus equivalence]
\label{thm:subcritical-locus-equivalence}
The standard logarithmic $abc$ conjecture is equivalent to
\eqref{eq:subcritical-locus-property}.
\end{theorem}

\begin{proof}
Assume $abc$, fix $0\leq\mu<1$, and set
$\eta=(1-\mu)/2$.  Then
\[
 d=1-(1+\eta)\mu=\frac{(1-\mu)(2-\mu)}2>0.
\]
If $C_\eta$ is the $abc$ constant and $R(P)\leq\mu H(P)$, then
\[
 dH(P)\leq C_\eta,
\]
so $B_\mu=C_\eta/d$ works.

Conversely, assume \eqref{eq:subcritical-locus-property}.  Given
$\eps>0$, put $\mu=(1+\eps)^{-1}$ and let $B_\mu$ be its bound.  If
$R(P)\leq\mu H(P)$, then
\[
 H(P)\leq B_\mu\leq(1+\eps)R(P)+\max\{B_\mu,0\}.
\]
If the displayed premise fails, then
$H(P)<(1+\eps)R(P)$.  Thus the single constant
$C_\eps=\max\{B_\mu,0\}$ works for every $P$.
\end{proof}

This is a proved equivalence, rather than a replacement definition.  It
shows that density zero, an exceptional-set estimate, or a positive power
saving cannot by itself prove $abc$: every point on every fixed
subcritical slope must be bounded.

\subsection{What counting would have to accomplish}

The logical gap persists even under very rapid separation.  Put
\[
                         X_n=2^{2^n}.
\]
Then $X_{n+1}=X_n^2$, but the sequence is infinite.  If
$A(X)=\#\{n:X_n\leq X\}$, then for every integer $X\geq2$,
\begin{equation}\label{eq:sparse-square-count}
                              A(X)^2\leq X.
\end{equation}
Indeed, two base-two logarithms give
$A(X)\leq 1+\lfloor\log_2\log_2X\rfloor$.  The elementary induction
$(N+1)^2\leq2^{2^N}$ completes \eqref{eq:sparse-square-count}.  Thus an
actual exponent-$1/2$ count and an exact square-gap law are compatible with
infinitely many objects.  This is a counterexample to that proposed counting
inference, not to $abc$.

Two upgrades do reach pointwise boundedness.  If the number of objects in
the $k$-th shell is at most $C2^{-\delta k}$ with $\delta>0$, then the
integer shell cardinality is eventually strictly below one and hence zero.
There is also a useful amplification criterion.

\begin{proposition}[Single-source amplification]
\label{prop:single-source-amplification}
Suppose $C>0$, $\theta,\kappa,\beta\geq0$, and a target class satisfies
$|T(Y)|\leq CY^\theta$ for $Y\geq1$.  If every bad source of height
$X\geq1$ produces a set of distinct targets $F_X$ with
\[
 F_X\subseteq T(X^\kappa),\qquad |F_X|\geq X^\beta,
 \qquad \beta>\kappa\theta,
\]
then every such source satisfies
\[
                  X\leq C^{1/(\beta-\kappa\theta)}.
\]
\end{proposition}

\begin{proof}
For each individual source,
\[
 X^\beta\leq|F_X|\leq|T(X^\kappa)|
       \leq CX^{\kappa\theta}.
\]
Division and the positive exponent $\beta-\kappa\theta$ give the result.
No overlap estimate between fibres of different sources is used.
\end{proof}

For the Browning--Verzobio bound in
Section~\ref{sec:three-prime-campana}, the resulting design inequality is
\[
 \beta>\kappa\left(\frac1p+\frac1q-\eta_{u,v}(r)+\eps\right).
\]
Their present positive exponent reaches neither an empty-shell bound nor
this amplification inequality without a new source-to-target construction.

\subsection{Residual-kernel escape and polynomial rigidity}

For $m\geq1$ and $n>0$, define
\[
 \kappa_m(n)=\prod_{\ell\mid n}\ell^{v_\ell(n)\bmod m},
 \qquad
 \rho_m(n)=\prod_{\ell\mid n}
       \ell^{\lfloor v_\ell(n)/m\rfloor}.
\]
Primewise Euclidean division gives the unique decomposition
\begin{equation}\label{eq:residual-kernel-decomposition}
                         n=\kappa_m(n)\rho_m(n)^m,
\end{equation}
with $\kappa_m(n)$ $m$-th-power-free.

\begin{proposition}[Kernel escape]
\label{prop:kernel-escape}
Fix $p,q,r\geq2$ with $1/p+1/q+1/r<1$.  In every infinite family of
distinct primitive points whose coordinates are respectively
$p$-, $q$-, and $r$-full, the triples
\[
       \bigl(\kappa_p(a),\kappa_q(b),\kappa_r(c)\bigr)
\]
take infinitely many values.
\end{proposition}

\begin{proof}
For a fixed residual triple $(A,B,C)$,
\eqref{eq:residual-kernel-decomposition} converts $a+b=c$ into
\[
                         Ax^p+By^q=Cz^r.
\]
Primitivity of $(a,b,c)$ implies $\gcd(x,y,z)=1$.  Darmon and
Granville~\cite[Theorem~2]{DarmonGranville1995} prove that this fixed
equation has only finitely many proper integer solutions in the strict
reciprocal range.  A finite residual packet would therefore give a finite
union of finite sets.
\end{proof}

The proposition rules out a finite list of fixed generalized-Fermat
coefficient triples.  It does not rule out a fixed Pell or elliptic source
curve whose coordinate values induce infinitely many residual kernels; the
Pell gate below has exactly this form.

There is a separate obstruction to putting all multiplicities into one
polynomial identity.

\begin{proposition}[Mason--Stothers obstruction]
\label{prop:polynomial-full-obstruction}
Let nonzero, pairwise coprime polynomials $A,B,C$ over a
characteristic-zero field satisfy $A+B=C$.  If they are respectively
polynomial-$p$-, $q$-, and $r$-full and at least one is nonconstant, then
\[
                         \frac1p+\frac1q+\frac1r>1.
\]
\end{proposition}

\begin{proof}
Put $D=\max\{\deg A,\deg B,\deg C\}>0$.  The polynomial $abc$ theorem of
Stothers~\cite{Stothers1981} gives
\[
 D+1\leq\deg\rad(ABC).
\]
Fullness and pairwise coprimality give
\[
 \deg\rad(ABC)
 \leq\frac{\deg A}{p}+\frac{\deg B}{q}+\frac{\deg C}{r}
 \leq D\left(\frac1p+\frac1q+\frac1r\right).
\]
The reciprocal sum cannot be at most one.
\end{proof}

This excludes polynomial-full identities.  It does not exclude a sparse set
of specializations at which polynomials with simple algebraic factors take
full integer values.

\subsection{A strict conditional Pell route}

The classical consecutive-powerful-number Pell orbit studied by
Walker~\cite{Walker1976}
\[
 x_0=1,\quad y_0=0,\qquad
 x_{n+1}=3x_n+8y_n,\quad y_{n+1}=x_n+3y_n
\]
satisfies
\begin{equation}\label{eq:pell-eight}
                     1+8y_n^2=x_n^2.
\end{equation}
Its positive coordinates are strictly increasing.  Hence it gives an
unbounded primitive family $(1,8y_n^2,x_n^2)$ in which both nonunit
coordinates are always squarefull.

If $y$ is itself squarefull, then $8y^2$ is cubefull: the prime $2$ occurs
to exponent at least three, and every odd prime occurs in $y^2$ to exponent
at least four.  The unit is $7$-full and $x^2$ is $2$-full, so
\[
                         (1,8y^2,x^2)
\]
has the strict signature $(7,3,2)$, with
\[
                       \frac17+\frac13+\frac12=\frac{41}{42}<1.
\]

\begin{theorem}[Pell-squarefull disproof gate]
\label{thm:pell-squarefull-disproof}
If an unbounded subsequence of the positive Pell roots $y_n$ is squarefull,
then the standard $abc$ conjecture is false.
\end{theorem}

\begin{proof}
The corresponding points from \eqref{eq:pell-eight} are positive and
primitive, have the fixed strict mixed-full signature $(7,3,2)$, and have
unbounded height.  The strict mixed-full disproof gate from
Section~\ref{sec:three-prime-campana} applies to the unchanged standard
target.
\end{proof}

The actual radical slope is stronger than $41/42$.  Squarefullness gives
$\rad(y)^2\leq y$; also $\rad(x)\leq x$ and $y<x$.  Therefore
\[
 \rad(8y^2x^2)^2
 \leq4\rad(y)^2\rad(x)^2
 \leq4yx^2<4x^3.
\]
With $H=\log(x^2)$ and $R=\log\rad(8y^2x^2)$ this yields
\begin{equation}\label{eq:pell-three-quarter-slope}
                         R\leq\frac34H+\log2.
\end{equation}
The unresolved premise is squarefullness of infinitely many $y_n$.
Here $\alpha=3+\sqrt8$ and $\beta=3-\sqrt8$ are algebraic integers with
$\alpha+\beta=6$, $\alpha\beta=1$, and coprime nonzero integral sum and
product; moreover $\alpha/\beta=\alpha^2>1$ is not a root of unity and
$y_n=(\alpha^n-\beta^n)/(\alpha-\beta)$.  Thus this is a nondegenerate Lucas
pair in the sense required by Bilu--Hanrot--Voutier.  Their Theorem~1.4 says
that every Lucas or Lehmer term of index greater than $30$ has a primitive
divisor~\cite{BHV2001}.  Their theorem does not assert that this divisor has
valuation one, so it alone does not negate the premise.  The first
twelve terms were factored only to check the recurrence and formulas; only
$y_1=1$ is squarefull among them, and no asymptotic conclusion is drawn.

The Lean companion proves the Pell norm recurrence and unboundedness, the
squarefull-to-cubefull transfer, the exact $(7,3,2)$ slope, and the conditional
implication to the original \texttt{ABCConjecture}.  It does not assume or prove an
unbounded squarefull subsequence.  The Darmon--Granville, Mason--Stothers,
three-quarter radical, and Lucas--Wieferich arguments remain paper proofs.
